\section{Discussion}
\label{sec:discussion}

\subsection{The Crossover Phenomenon}
\label{subsec:crossover_discussion}

Our experimental results reveal a sharp, predictable crossover boundary that is governed by a simple relationship between the memory budget $M$ and the per-element cost of each data structure.

\textbf{Hash Table capacity ceiling.} The maximum number of elements a Hash Table can store within a budget $M$ is:
\begin{equation}
    N_{\max}^{\text{HT}} = \left\lfloor \frac{M}{C_{\text{HT}}} \right\rfloor \approx \frac{M}{100~\text{bytes}}
    \label{eq:ht_capacity}
\end{equation}
where $C_{\text{HT}} \approx 100$ bytes includes the node overhead (32 bytes), string object (52 bytes), and table reference (8 bytes) plus alignment padding. Our measured values confirm this: at $M = 64$~MB, the Hash Table accommodates up to $N = 500{,}000$ ($\approx 47$~MB) but fails at $N = 1{,}000{,}000$ ($\approx 95$~MB).

\textbf{Bloom Filter capacity ceiling.} By contrast:
\begin{equation}
    N_{\max}^{\text{BF}} = \left\lfloor \frac{M}{C_{\text{BF}}} \right\rfloor \approx \frac{M}{1.2~\text{bytes}}
    \label{eq:bf_capacity}
\end{equation}
For a 64~MB budget, this yields $N_{\max}^{\text{BF}} \approx 53{,}000{,}000$---over $80\times$ the Hash Table's capacity. In our experiments, the Bloom Filter at $N = 2{,}000{,}000$ uses only 2.3~MB, leaving the vast majority of the budget unused.

\textbf{Pre-OOM degradation.} A critical finding is that the Hash Table does not fail gracefully at the capacity ceiling. Instead, the JVM garbage collector begins aggressive stop-the-world collection cycles as heap utilization exceeds ${\approx}70\%$ of the budget. This manifests as throughput degradation \emph{before} the actual \texttt{OutOfMemoryError}. In production systems, this GC-induced latency spike is often more damaging than a clean crash, as it causes cascading timeouts~\citep{jvmspec2023}.

\subsection{Interpreting the Compression Ratio}
\label{subsec:compression}

The $83.5\times$ compression ratio is a direct consequence of the fundamental design difference:
\begin{itemize}[nosep]
    \item The Hash Table \emph{stores} each key, requiring memory proportional to the key's content.
    \item The Bloom Filter \emph{encodes} membership as a bit pattern, requiring only ${\approx}9.6$ bits per element regardless of key length.
\end{itemize}

This ratio remains constant because both structures have \emph{linear} memory complexity ($O(n)$), and our datasets use a uniform key format. In practice, the ratio would increase for longer keys (e.g., URLs, email addresses), making the Bloom Filter even more advantageous.

\subsection{FPR as an Acceptable Trade-off}
\label{subsec:fpr_tradeoff}

Our empirical FPR of $0.84\%$--$1.38\%$ closely tracks the theoretical target of $1\%$. The deviations are consistent with binomial sampling variance: for 5{,}000 queries with $p = 0.01$, the standard deviation is $\sigma = \sqrt{np(1-p)} \approx 7$, yielding a 95\% confidence interval of approximately $0.72\%$--$1.28\%$.

For many real-world applications, a 1\% FPR is an acceptable cost:
\begin{itemize}[nosep]
    \item \textbf{CDN/Cache filtering}: A 1\% FPR means that out of 1{,}000 negative queries, only ${\approx}10$ trigger an unnecessary (but harmless) cache lookup.
    \item \textbf{Database pre-filtering}: The Bloom Filter blocks 99\% of unnecessary disk I/O operations, with only 1\% passing through as false positives to be resolved by a secondary exact check.
    \item \textbf{Network deduplication}: In systems processing millions of packets per second, a 1\% FPR translates to negligible overhead compared to the cost of storing every packet hash in a full Hash Table.
\end{itemize}

\subsection{Practical Decision Framework}
\label{subsec:framework}

Based on our findings, we propose the following decision guideline:

\begin{table}[h]
\centering
\caption{Practical decision framework for choosing between Hash Table and Bloom Filter.}
\label{tab:decision}
\begin{tabular}{@{}p{3.8cm}p{5.5cm}p{5.5cm}@{}}
\toprule
\textbf{Criterion} & \textbf{Choose Hash Table} & \textbf{Choose Bloom Filter} \\
\midrule
Accuracy requirement & Must be 100\% exact (zero tolerance for false positives) & Can tolerate ${\leq}1\%$ false positive rate \\
\addlinespace
Memory availability  & $N < N_{\max}^{\text{HT}} = M / 100$ & $N > N_{\max}^{\text{HT}}$ or memory is scarce \\
\addlinespace
Value retrieval      & Need to retrieve stored values (key--value mapping) & Only need existence check (yes/no) \\
\addlinespace
Deletion support     & Support required (can remove keys) & Not needed (standard BF cannot delete) \\
\addlinespace
Use case             & Small--medium sets, exact databases & Large-scale filtering, pre-screening, caching \\
\bottomrule
\end{tabular}
\end{table}

\subsection{Threats to Validity}
\label{subsec:threats}

Several factors may affect the generalizability of our results:

\begin{enumerate}[nosep, label=(\roman*)]
    \item \textbf{JVM overhead}: The JVM's own memory consumption (class metadata, thread stacks, JIT compilation buffers) reduces the effective heap available for data structures. This causes the observed OOM thresholds to differ from the pure theoretical $N_{\max}$.
    \item \textbf{Dataset uniformity}: Our datasets use uniformly distributed random usernames. Highly skewed or adversarial key distributions may alter Hash Table collision behavior.
    \item \textbf{Single-threaded execution}: Our benchmarks are single-threaded. Concurrent access patterns may introduce contention effects, particularly for Hash Table resizing operations.
    \item \textbf{Throughput variance}: JIT warm-up, GC pauses, and OS scheduling introduce noise in throughput measurements. We mitigate this with warm-up iterations but do not perform multiple trial averaging.
\end{enumerate}
